\section{Learning Objectives}

After completing this chapter, readers will be able to:
\begin{itemize}
    \item Design data pipeline architectures for batch, streaming, and real-time ML workloads
    \item Choose between ETL and ELT patterns based on ML requirements
    \item Implement comprehensive data validation using Great Expectations
    \item Build and operate feature stores for ML feature management
    \item Handle schema evolution and ensure backward compatibility
    \item Use Apache Beam for portable data processing pipelines
    \item Integrate Kafka for streaming data ingestion
    \item Version datasets and track data lineage with DVC
    \item Monitor data quality and detect data drift
    \item Design fault-tolerant, scalable data pipelines
\end{itemize}

\section{Introduction: Data Pipelines as the Foundation}

Data pipelines form the critical foundation of any production ML system. While models receive the majority of attention, the reality is that 80\% of ML engineering effort goes into data preparation, quality assurance, and pipeline maintenance \citep{cognilytica2020}.

A robust data pipeline must address several challenges simultaneously:

\begin{itemize}
    \item \textbf{Scale:} Processing gigabytes to petabytes of data efficiently
    \item \textbf{Reliability:} Ensuring consistent, accurate data delivery
    \item \textbf{Latency:} Meeting timing requirements from batch to real-time
    \item \textbf{Quality:} Validating data and detecting anomalies
    \item \textbf{Versioning:} Tracking data evolution and enabling reproducibility
    \item \textbf{Governance:} Ensuring privacy, security, and compliance
\end{itemize}

Unlike traditional software systems where data structures are relatively stable, ML pipelines must handle evolving data distributions, schema changes, and quality degradation—all while maintaining model performance.

\subsection{The Cost of Poor Data Pipelines}

Poor data pipeline design leads to catastrophic failures:

\begin{example}[Data Pipeline Failure]
A recommendation system at a major e-commerce company experienced a 40\% drop in click-through rate overnight. Investigation revealed:

\begin{itemize}
    \item The upstream data pipeline silently failed to update user preferences
    \item Models trained on stale data (3 weeks old) instead of daily updates
    \item No data freshness validation caught the issue
    \item Impact: \$2M in lost revenue over 5 days
\end{itemize}

The failure stemmed from lack of data validation and monitoring, not model issues.
\end{example}

\section{Data Ingestion Patterns}

Data ingestion patterns fundamentally differ in their latency, complexity, and cost characteristics. Understanding these trade-offs is essential for appropriate pattern selection.

\subsection{Batch Ingestion}

Batch ingestion processes data in large chunks at scheduled intervals—hourly, daily, or weekly. This remains the most common pattern for ML training pipelines.

\subsubsection{Characteristics}

\begin{itemize}
    \item \textbf{Latency:} Minutes to hours between data generation and availability
    \item \textbf{Throughput:} High—optimized for large-volume processing
    \item \textbf{Cost:} Low—resource utilization is predictable and optimizable
    \item \textbf{Complexity:} Low—simpler to implement and debug
    \item \textbf{Use Cases:} Training data preparation, feature engineering, historical analysis
\end{itemize}

\subsubsection{Architecture}

\begin{figure}[h]
\centering
\begin{tikzpicture}[
    node distance=2cm,
    box/.style={rectangle, draw, fill=blue!20, text width=2.5cm, text centered, minimum height=1cm},
    storage/.style={cylinder, draw, fill=green!20, text width=2cm, text centered, minimum height=1cm, aspect=0.3},
    arrow/.style={->, >=stealth, thick}
]
    % Nodes
    \node[box] (source) {Data Sources};
    \node[storage, right of=source, xshift=2cm] (landing) {Landing Zone};
    \node[box, right of=landing, xshift=2.5cm] (etl) {ETL Pipeline};
    \node[storage, right of=etl, xshift=2cm] (warehouse) {Data Warehouse};
    \node[box, below of=warehouse, yshift=-1cm] (feature) {Feature Engineering};
    \node[storage, left of=feature, xshift=-2cm] (feature_store) {Feature Store};

    % Scheduler
    \node[box, above of=etl, yshift=1cm] (scheduler) {Scheduler (Airflow)};

    % Arrows
    \draw[arrow] (source) -- (landing);
    \draw[arrow] (landing) -- (etl);
    \draw[arrow] (etl) -- (warehouse);
    \draw[arrow] (warehouse) -- (feature);
    \draw[arrow] (feature) -- (feature_store);
    \draw[arrow, dashed] (scheduler) -- (etl);

    % Labels
    \node[above, text width=2cm, align=center] at ($(source)!0.5!(landing)$) {\small Daily Batch};
    \node[above, text width=2cm, align=center] at ($(landing)!0.5!(etl)$) {\small Validate \& Transform};
\end{tikzpicture}
\caption{Batch Data Ingestion Architecture}
\label{fig:batch_architecture}
\end{figure}

\subsubsection{Implementation with Apache Beam}

Apache Beam provides a unified programming model for batch and streaming data processing:

\begin{lstlisting}[style=python, caption=Batch Data Pipeline with Apache Beam]
import apache_beam as beam
from apache_beam.options.pipeline_options import PipelineOptions
from apache_beam.io import ReadFromText, WriteToBigQuery
from typing import Dict, Any
import logging
from datetime import datetime

class ValidateData(beam.DoFn):
    """Validate incoming data records"""

    def __init__(self, schema: Dict[str, type]):
        self.schema = schema
        self.valid_count = beam.metrics.Metrics.counter('main', 'valid_records')
        self.invalid_count = beam.metrics.Metrics.counter('main', 'invalid_records')

    def process(self, element: str):
        """Validate and parse CSV record"""
        try:
            # Parse CSV
            fields = element.split(',')
            record = {
                'user_id': fields[0],
                'timestamp': fields[1],
                'event_type': fields[2],
                'value': float(fields[3])
            }

            # Validate schema
            if not self._is_valid(record):
                self.invalid_count.inc()
                logging.warning(f"Invalid record: {record}")
                yield beam.pvalue.TaggedOutput('invalid', record)
            else:
                self.valid_count.inc()
                yield record

        except Exception as e:
            self.invalid_count.inc()
            logging.error(f"Parse error: {e}")
            yield beam.pvalue.TaggedOutput('invalid', element)

    def _is_valid(self, record: Dict[str, Any]) -> bool:
        """Validate record against schema"""
        # Check required fields
        required_fields = ['user_id', 'timestamp', 'event_type', 'value']
        if not all(field in record for field in required_fields):
            return False

        # Check data types
        if not isinstance(record['value'], (int, float)):
            return False

        # Check value ranges
        if record['value'] < 0:
            return False

        # Check timestamp format
        try:
            datetime.fromisoformat(record['timestamp'])
        except ValueError:
            return False

        return True

class EnrichData(beam.DoFn):
    """Enrich records with additional features"""

    def process(self, element: Dict[str, Any]):
        """Add derived features"""
        element['hour'] = datetime.fromisoformat(
            element['timestamp']
        ).hour

        element['date'] = datetime.fromisoformat(
            element['timestamp']
        ).date().isoformat()

        element['day_of_week'] = datetime.fromisoformat(
            element['timestamp']
        ).weekday()

        yield element

class AggregateMetrics(beam.DoFn):
    """Aggregate metrics by user"""

    def process(self, element: tuple):
        """Calculate user-level aggregations"""
        user_id, records = element
        records_list = list(records)

        aggregated = {
            'user_id': user_id,
            'total_events': len(records_list),
            'total_value': sum(r['value'] for r in records_list),
            'avg_value': sum(r['value'] for r in records_list) / len(records_list),
            'unique_event_types': len(set(r['event_type'] for r in records_list)),
            'date': records_list[0]['date']
        }

        yield aggregated

def run_batch_pipeline(
    input_path: str,
    output_table: str,
    pipeline_options: PipelineOptions
):
    """Execute batch data pipeline"""

    schema = {
        'user_id': str,
        'timestamp': str,
        'event_type': str,
        'value': float
    }

    with beam.Pipeline(options=pipeline_options) as pipeline:

        # Read data
        raw_data = (
            pipeline
            | 'ReadFromGCS' >> ReadFromText(input_path, skip_header_lines=1)
        )

        # Validate and parse
        validated = (
            raw_data
            | 'ValidateData' >> beam.ParDo(ValidateData(schema)).with_outputs(
                'invalid',
                main='valid'
            )
        )

        # Process valid records
        enriched = (
            validated.valid
            | 'EnrichData' >> beam.ParDo(EnrichData())
        )

        # Aggregate by user
        aggregated = (
            enriched
            | 'KeyByUser' >> beam.Map(lambda x: (x['user_id'], x))
            | 'GroupByUser' >> beam.GroupByKey()
            | 'Aggregate' >> beam.ParDo(AggregateMetrics())
        )

        # Write to BigQuery
        _ = (
            aggregated
            | 'WriteToBigQuery' >> WriteToBigQuery(
                output_table,
                schema='user_id:STRING,total_events:INTEGER,total_value:FLOAT,'
                       'avg_value:FLOAT,unique_event_types:INTEGER,date:DATE',
                write_disposition=beam.io.BigQueryDisposition.WRITE_APPEND,
                create_disposition=beam.io.BigQueryDisposition.CREATE_IF_NEEDED
            )
        )

        # Write invalid records to dead letter queue
        _ = (
            validated.invalid
            | 'WriteInvalid' >> beam.io.WriteToText(
                f'{output_table}_invalid',
                file_name_suffix='.json'
            )
        )

# Run pipeline
if __name__ == '__main__':
    import argparse

    parser = argparse.ArgumentParser()
    parser.add_argument('--input', required=True, help='Input file pattern')
    parser.add_argument('--output', required=True, help='Output BigQuery table')
    parser.add_argument('--runner', default='DirectRunner',
                       help='Pipeline runner (DirectRunner, DataflowRunner)')

    args, beam_args = parser.parse_known_args()

    pipeline_options = PipelineOptions(
        beam_args,
        runner=args.runner,
        project='my-project',
        temp_location='gs://my-bucket/temp',
        region='us-central1'
    )

    run_batch_pipeline(args.input, args.output, pipeline_options)
\end{lstlisting}

\subsection{Streaming Ingestion}

Streaming ingestion processes data continuously as it arrives, enabling near-real-time analytics and model updates.

\subsubsection{Characteristics}

\begin{itemize}
    \item \textbf{Latency:} Seconds to minutes
    \item \textbf{Throughput:} Medium to high—depends on parallelization
    \item \textbf{Cost:} Medium to high—requires continuous resource allocation
    \item \textbf{Complexity:} High—windowing, state management, fault tolerance
    \item \textbf{Use Cases:} Real-time features, continuous model monitoring, live dashboards
\end{itemize}

\subsubsection{Architecture with Kafka}

\begin{figure}[h]
\centering
\begin{tikzpicture}[
    node distance=2cm,
    box/.style={rectangle, draw, fill=blue!20, text width=2.5cm, text centered, minimum height=1cm},
    kafka/.style={rectangle, draw, fill=orange!30, text width=2.5cm, text centered, minimum height=1cm},
    storage/.style={cylinder, draw, fill=green!20, text width=2cm, text centered, minimum height=1cm, aspect=0.3},
    arrow/.style={->, >=stealth, thick}
]
    % Producer side
    \node[box] (producer) {Event Producers};
    \node[kafka, right of=producer, xshift=2cm] (kafka) {Kafka Topics};

    % Consumer side
    \node[box, below of=kafka, yshift=-2cm] (consumer1) {Stream Processor};
    \node[box, right of=consumer1, xshift=3cm] (consumer2) {ML Feature Extractor};
    \node[box, left of=consumer1, xshift=-3cm] (consumer3) {Analytics};

    % Storage
    \node[storage, below of=consumer2, yshift=-1.5cm] (feature_store) {Feature Store};
    \node[storage, below of=consumer1, yshift=-1.5cm] (warehouse) {Data Warehouse};

    % Arrows
    \draw[arrow] (producer) -- (kafka);
    \draw[arrow] (kafka) -- (consumer1);
    \draw[arrow] (kafka) -- (consumer2);
    \draw[arrow] (kafka) -- (consumer3);
    \draw[arrow] (consumer2) -- (feature_store);
    \draw[arrow] (consumer1) -- (warehouse);

    % Labels
    \node[above right, text width=2cm] at (producer.north east) {\tiny Real-time Events};
    \node[right, text width=2.5cm] at (kafka.east) {\tiny Durable, Partitioned Log};
\end{tikzpicture}
\caption{Streaming Data Architecture with Kafka}
\label{fig:streaming_architecture}
\end{figure}

\subsubsection{Implementation with Kafka and Flink}

\begin{lstlisting}[style=python, caption=Streaming Pipeline with Kafka]
from kafka import KafkaConsumer, KafkaProducer
from typing import Dict, Any, Optional
import json
import logging
from datetime import datetime
from collections import defaultdict
import threading
import time

logger = logging.getLogger(__name__)

class StreamingFeatureExtractor:
    """Extract features from streaming data"""

    def __init__(
        self,
        kafka_bootstrap_servers: str,
        input_topic: str,
        output_topic: str,
        group_id: str
    ):
        self.bootstrap_servers = kafka_bootstrap_servers
        self.input_topic = input_topic
        self.output_topic = output_topic
        self.group_id = group_id

        # State for windowed aggregations
        self.window_state = defaultdict(lambda: defaultdict(list))
        self.window_duration = 300  # 5 minutes

        # Initialize Kafka consumer
        self.consumer = KafkaConsumer(
            input_topic,
            bootstrap_servers=kafka_bootstrap_servers,
            group_id=group_id,
            value_deserializer=lambda m: json.loads(m.decode('utf-8')),
            auto_offset_reset='latest',
            enable_auto_commit=True,
            max_poll_records=500
        )

        # Initialize Kafka producer
        self.producer = KafkaProducer(
            bootstrap_servers=kafka_bootstrap_servers,
            value_serializer=lambda v: json.dumps(v).encode('utf-8'),
            acks='all',
            retries=3
        )

        # Start window cleanup thread
        self.running = True
        self.cleanup_thread = threading.Thread(
            target=self._cleanup_old_windows,
            daemon=True
        )
        self.cleanup_thread.start()

    def process_event(self, event: Dict[str, Any]) -> Optional[Dict[str, Any]]:
        """Process individual event and extract features"""
        try:
            # Validate event
            if not self._validate_event(event):
                logger.warning(f"Invalid event: {event}")
                return None

            # Extract base features
            features = {
                'user_id': event['user_id'],
                'timestamp': event['timestamp'],
                'event_type': event['event_type'],
                'value': event['value'],
            }

            # Add time-based features
            dt = datetime.fromisoformat(event['timestamp'])
            features['hour'] = dt.hour
            features['day_of_week'] = dt.weekday()
            features['is_weekend'] = 1 if dt.weekday() >= 5 else 0

            # Update window state
            window_key = self._get_window_key(dt)
            user_id = event['user_id']
            self.window_state[user_id][window_key].append(event['value'])

            # Calculate windowed features
            window_features = self._calculate_window_features(user_id, window_key)
            features.update(window_features)

            return features

        except Exception as e:
            logger.error(f"Error processing event: {e}", exc_info=True)
            return None

    def _validate_event(self, event: Dict[str, Any]) -> bool:
        """Validate event schema and values"""
        required_fields = ['user_id', 'timestamp', 'event_type', 'value']

        if not all(field in event for field in required_fields):
            return False

        # Validate value range
        if not isinstance(event['value'], (int, float)) or event['value'] < 0:
            return False

        # Validate timestamp
        try:
            datetime.fromisoformat(event['timestamp'])
        except ValueError:
            return False

        return True

    def _get_window_key(self, dt: datetime) -> int:
        """Get window key for timestamp"""
        return int(dt.timestamp() // self.window_duration)

    def _calculate_window_features(
        self,
        user_id: str,
        current_window: int
    ) -> Dict[str, float]:
        """Calculate features from windowed data"""

        # Get data from current and previous windows
        windows_to_check = [current_window, current_window - 1]
        all_values = []

        for window_key in windows_to_check:
            if window_key in self.window_state[user_id]:
                all_values.extend(self.window_state[user_id][window_key])

        if not all_values:
            return {
                'window_count': 0,
                'window_sum': 0.0,
                'window_avg': 0.0,
                'window_max': 0.0
            }

        return {
            'window_count': len(all_values),
            'window_sum': sum(all_values),
            'window_avg': sum(all_values) / len(all_values),
            'window_max': max(all_values)
        }

    def _cleanup_old_windows(self):
        """Remove old window data to prevent memory growth"""
        while self.running:
            try:
                current_time = datetime.now().timestamp()
                current_window = int(current_time // self.window_duration)

                # Keep only last 3 windows
                cutoff_window = current_window - 3

                for user_id in list(self.window_state.keys()):
                    windows_to_remove = [
                        w for w in self.window_state[user_id].keys()
                        if w < cutoff_window
                    ]

                    for window in windows_to_remove:
                        del self.window_state[user_id][window]

                    # Remove user if no windows left
                    if not self.window_state[user_id]:
                        del self.window_state[user_id]

                time.sleep(60)  # Run every minute

            except Exception as e:
                logger.error(f"Error in cleanup: {e}", exc_info=True)

    def run(self):
        """Main processing loop"""
        logger.info(f"Starting streaming processor for topic {self.input_topic}")

        try:
            for message in self.consumer:
                event = message.value

                # Process event
                features = self.process_event(event)

                if features:
                    # Send to output topic
                    self.producer.send(
                        self.output_topic,
                        value=features,
                        key=features['user_id'].encode('utf-8')
                    )

                    # Flush periodically
                    if message.offset % 100 == 0:
                        self.producer.flush()

        except KeyboardInterrupt:
            logger.info("Shutting down...")

        finally:
            self.shutdown()

    def shutdown(self):
        """Clean shutdown"""
        self.running = False
        self.consumer.close()
        self.producer.flush()
        self.producer.close()
        logger.info("Shutdown complete")

# Usage
if __name__ == '__main__':
    import argparse

    parser = argparse.ArgumentParser()
    parser.add_argument('--bootstrap-servers', required=True)
    parser.add_argument('--input-topic', required=True)
    parser.add_argument('--output-topic', required=True)
    parser.add_argument('--group-id', required=True)

    args = parser.parse_args()

    processor = StreamingFeatureExtractor(
        kafka_bootstrap_servers=args.bootstrap_servers,
        input_topic=args.input_topic,
        output_topic=args.output_topic,
        group_id=args.group_id
    )

    processor.run()
\end{lstlisting}

\subsection{Hybrid Patterns: Lambda Architecture}

Many production ML systems require both batch and streaming processing. Lambda architecture addresses this by maintaining separate batch and streaming pipelines.

\begin{figure}[h]
\centering
\begin{tikzpicture}[
    node distance=1.5cm,
    box/.style={rectangle, draw, fill=blue!20, text width=2.3cm, text centered, minimum height=0.8cm, font=\small},
    storage/.style={cylinder, draw, fill=green!20, text width=1.8cm, text centered, minimum height=0.8cm, aspect=0.3, font=\small},
    arrow/.style={->, >=stealth, thick}
]
    % Data source
    \node[storage] (source) {Data Source};

    % Batch layer
    \node[box, below left of=source, xshift=-2cm, yshift=-1.5cm] (batch) {Batch Layer};
    \node[storage, below of=batch, yshift=-0.5cm] (batch_views) {Batch Views};

    % Speed layer
    \node[box, below right of=source, xshift=2cm, yshift=-1.5cm] (speed) {Speed Layer};
    \node[storage, below of=speed, yshift=-0.5cm] (realtime_views) {RT Views};

    % Serving layer
    \node[box, below of=source, yshift=-5cm] (serving) {Serving Layer};

    % Arrows
    \draw[arrow] (source) -| (batch);
    \draw[arrow] (source) -| (speed);
    \draw[arrow] (batch) -- (batch_views);
    \draw[arrow] (speed) -- (realtime_views);
    \draw[arrow] (batch_views) -- (serving);
    \draw[arrow] (realtime_views) -- (serving);

    % Labels
    \node[left, text width=2cm, font=\tiny] at (batch.west) {Complete, accurate};
    \node[right, text width=2cm, font=\tiny] at (speed.east) {Low latency, approximate};
\end{tikzpicture}
\caption{Lambda Architecture for ML}
\label{fig:lambda_architecture}
\end{figure}

\section{ETL vs ELT for ML Workloads}

The choice between ETL (Extract, Transform, Load) and ELT (Extract, Load, Transform) significantly impacts ML pipeline design.

\subsection{ETL: Traditional Approach}

ETL transforms data before loading into the target system.

\textbf{Advantages for ML:}
\begin{itemize}
    \item Cleaner data in warehouse—only curated data persisted
    \item Lower storage costs—unnecessary data filtered early
    \item Preprocessing embedded in pipeline
    \item Suitable for well-defined transformations
\end{itemize}

\textbf{Disadvantages for ML:}
\begin{itemize}
    \item Loses raw data—hard to reprocess with new logic
    \item Rigid—changing transformations requires pipeline updates
    \item Limited exploratory analysis—transformed data may not suit all use cases
    \item Difficult to debug data issues—original data not preserved
\end{itemize}

\subsection{ELT: Modern Approach}

ELT loads raw data first, then transforms within the target system.

\textbf{Advantages for ML:}
\begin{itemize}
    \item Preserves raw data—enables reprocessing and debugging
    \item Flexible—can apply different transformations for different models
    \item Supports exploration—data scientists can experiment freely
    \item Reproducible—transformations versioned with models
    \item Leverages warehouse compute—use BigQuery, Snowflake, etc.
\end{itemize}

\textbf{Disadvantages for ML:}
\begin{itemize}
    \item Higher storage costs—keeping raw and transformed data
    \item Requires powerful target system
    \item Potential for data duplication
    \item More complex data governance
\end{itemize}

\begin{bestpractice}
For ML workloads, prefer ELT:
\begin{itemize}
    \item Store raw data in data lake (S3, GCS)
    \item Load into data warehouse (BigQuery, Snowflake)
    \item Apply transformations using SQL or Spark
    \item Version transformations with models
    \item Enables experimentation and reproducibility
\end{itemize}
\end{bestpractice}

\subsection{Hybrid Approach for ML}

Production ML systems often use a hybrid approach:

\begin{enumerate}
    \item \textbf{Raw Data Lake:} Store all raw data (ELT)
    \item \textbf{Staging Layer:} Basic cleaning and standardization (ETL)
    \item \textbf{Feature Engineering:} Model-specific transformations (ELT)
    \item \textbf{Feature Store:} Curated, versioned features
\end{enumerate}

\section{Data Validation and Quality Assurance}

Data quality is the most critical—and often overlooked—aspect of ML pipelines. Poor data quality causes more production failures than model issues.

\subsection{Data Quality Dimensions}

\begin{table}[h]
\centering
\caption{Data Quality Dimensions for ML}
\label{tab:data_quality}
\begin{tabular}{@{}llp{6cm}@{}}
\toprule
\textbf{Dimension} & \textbf{Question} & \textbf{ML Impact} \\
\midrule
Accuracy & Is data correct? & Models learn incorrect patterns \\
Completeness & Are all required fields present? & Missing features degrade performance \\
Consistency & Is data consistent across sources? & Training-serving skew \\
Timeliness & Is data current? & Models become stale \\
Validity & Does data conform to rules? & Invalid inputs cause errors \\
Uniqueness & Are duplicates removed? & Biased training data \\
\bottomrule
\end{tabular}
\end{table}

\subsection{Data Validation with Great Expectations}

Great Expectations provides a framework for data validation, documentation, and profiling.

\begin{lstlisting}[style=python, caption=Comprehensive Data Validation with Great Expectations]
import great_expectations as ge
from great_expectations.core import ExpectationSuite
from great_expectations.data_context import DataContext
from typing import Dict, Any
import pandas as pd
import logging

logger = logging.getLogger(__name__)

class MLDataValidator:
    """Comprehensive data validation for ML pipelines"""

    def __init__(self, context_root_dir: str):
        """Initialize with Great Expectations context"""
        self.context = DataContext(context_root_dir)

    def create_ml_expectation_suite(
        self,
        suite_name: str,
        sample_df: pd.DataFrame
    ) -> ExpectationSuite:
        """Create comprehensive expectation suite for ML data"""

        # Create validator
        validator = self.context.get_validator(
            batch_request={
                'datasource_name': 'pandas',
                'data_asset_name': suite_name
            },
            expectation_suite_name=suite_name
        )

        # Table-level expectations
        validator.expect_table_row_count_to_be_between(
            min_value=1000,
            max_value=100000000,
            meta={'notes': 'Ensure sufficient data for training'}
        )

        validator.expect_table_column_count_to_equal(
            value=len(sample_df.columns),
            meta={'notes': 'Schema stability check'}
        )

        # Numeric column expectations
        numeric_cols = sample_df.select_dtypes(include=['float64', 'int64']).columns

        for col in numeric_cols:
            # Non-null
            validator.expect_column_values_to_not_be_null(
                column=col,
                mostly=0.95,  # Allow up to 5% nulls
                meta={'notes': f'{col} should be mostly complete'}
            )

            # Value range based on training data
            q1 = sample_df[col].quantile(0.01)
            q99 = sample_df[col].quantile(0.99)

            validator.expect_column_values_to_be_between(
                column=col,
                min_value=q1 * 0.5,  # Allow some deviation
                max_value=q99 * 1.5,
                mostly=0.95,
                meta={'notes': f'{col} range check against training distribution'}
            )

            # Check for excessive zeros (potential data quality issue)
            validator.expect_column_proportion_of_unique_values_to_be_between(
                column=col,
                min_value=0.01,
                max_value=1.0,
                meta={'notes': f'{col} should have reasonable diversity'}
            )

        # Categorical column expectations
        categorical_cols = sample_df.select_dtypes(include=['object']).columns

        for col in categorical_cols:
            # Non-null
            validator.expect_column_values_to_not_be_null(
                column=col,
                mostly=0.98
            )

            # Known categories (from training)
            known_values = sample_df[col].unique().tolist()

            validator.expect_column_values_to_be_in_set(
                column=col,
                value_set=known_values,
                mostly=0.90,  # Allow some new categories
                meta={'notes': f'{col} mostly known categories'}
            )

            # Cardinality check
            validator.expect_column_unique_value_count_to_be_between(
                column=col,
                min_value=1,
                max_value=len(known_values) * 1.2,  # Allow 20% new categories
                meta={'notes': f'{col} cardinality check'}
            )

        # Cross-column expectations
        # Example: start_date should be before end_date
        if 'start_date' in sample_df.columns and 'end_date' in sample_df.columns:
            validator.expect_column_pair_values_A_to_be_greater_than_B(
                column_A='end_date',
                column_B='start_date',
                or_equal=True,
                mostly=0.99
            )

        # Statistical distribution checks
        for col in numeric_cols:
            # Check if distribution has shifted significantly
            expected_mean = sample_df[col].mean()
            expected_std = sample_df[col].std()

            validator.expect_column_mean_to_be_between(
                column=col,
                min_value=expected_mean - 3 * expected_std,
                max_value=expected_mean + 3 * expected_std,
                meta={'notes': f'{col} mean should be stable'}
            )

        # Save suite
        validator.save_expectation_suite(discard_failed_expectations=False)

        return validator.get_expectation_suite()

    def validate_data(
        self,
        df: pd.DataFrame,
        suite_name: str,
        run_name: str = None
    ) -> Dict[str, Any]:
        """Validate data against expectation suite"""

        # Create batch
        batch = self.context.get_batch(
            batch_request={
                'datasource_name': 'pandas',
                'data_asset_name': suite_name,
                'data': df
            },
            expectation_suite_name=suite_name
        )

        # Run validation
        results = batch.validate()

        # Parse results
        validation_results = {
            'success': results.success,
            'evaluated_expectations': results.statistics['evaluated_expectations'],
            'successful_expectations': results.statistics['successful_expectations'],
            'unsuccessful_expectations': results.statistics['unsuccessful_expectations'],
            'success_percent': results.statistics['success_percent'],
            'failed_expectations': []
        }

        # Collect failed expectations
        for result in results.results:
            if not result.success:
                validation_results['failed_expectations'].append({
                    'expectation_type': result.expectation_config.expectation_type,
                    'kwargs': result.expectation_config.kwargs,
                    'result': result.result
                })

        # Log results
        if results.success:
            logger.info(f"Validation passed: {validation_results['success_percent']:.1f}%")
        else:
            logger.error(
                f"Validation failed: {len(validation_results['failed_expectations'])} "
                f"expectations failed"
            )
            for failed in validation_results['failed_expectations'][:5]:
                logger.error(f"  - {failed['expectation_type']}: {failed['kwargs']}")

        # Build and save data docs (optional)
        if run_name:
            self.context.build_data_docs()

        return validation_results

    def detect_data_drift(
        self,
        reference_df: pd.DataFrame,
        current_df: pd.DataFrame,
        threshold: float = 0.05
    ) -> Dict[str, Any]:
        """Detect data drift using statistical tests"""
        from scipy import stats

        drift_results = {}

        numeric_cols = reference_df.select_dtypes(
            include=['float64', 'int64']
        ).columns

        for col in numeric_cols:
            # Kolmogorov-Smirnov test
            statistic, p_value = stats.ks_2samp(
                reference_df[col].dropna(),
                current_df[col].dropna()
            )

            drift_detected = p_value < threshold

            drift_results[col] = {
                'drift_detected': drift_detected,
                'p_value': p_value,
                'ks_statistic': statistic,
                'ref_mean': reference_df[col].mean(),
                'current_mean': current_df[col].mean(),
                'mean_shift': (
                    current_df[col].mean() - reference_df[col].mean()
                ) / reference_df[col].std()
            }

            if drift_detected:
                logger.warning(
                    f"Drift detected in {col}: p-value={p_value:.4f}, "
                    f"mean shift={drift_results[col]['mean_shift']:.2f} std"
                )

        return drift_results

# Usage example
def validate_ml_data():
    """Example usage of ML data validator"""

    # Initialize validator
    validator = MLDataValidator(context_root_dir='./great_expectations')

    # Load training data (reference)
    train_df = pd.read_csv('data/train.csv')

    # Create expectation suite based on training data
    suite = validator.create_ml_expectation_suite(
        suite_name='ml_training_data',
        sample_df=train_df
    )

    # Load new data to validate
    new_df = pd.read_csv('data/new_batch.csv')

    # Validate
    results = validator.validate_data(
        df=new_df,
        suite_name='ml_training_data',
        run_name='daily_validation'
    )

    # Check for drift
    drift_results = validator.detect_data_drift(
        reference_df=train_df,
        current_df=new_df,
        threshold=0.05
    )

    # Decision logic
    if not results['success']:
        raise ValueError(
            f"Data validation failed: {results['unsuccessful_expectations']} "
            f"expectations failed"
        )

    drifted_features = [
        feat for feat, result in drift_results.items()
        if result['drift_detected']
    ]

    if len(drifted_features) > 5:
        logger.warning(
            f"Significant drift detected in {len(drifted_features)} features. "
            f"Consider retraining."
        )

    return results, drift_results

if __name__ == '__main__':
    results, drift = validate_ml_data()
\end{lstlisting}

This comprehensive validation implementation:
- Creates expectations automatically from training data
- Validates new data against these expectations
- Detects statistical drift using KS tests
- Provides detailed logging and reporting
- Enables data quality gates in pipelines

\section{Feature Stores: Centralized Feature Management}

Feature stores solve a critical challenge in ML systems: ensuring consistency between training and serving features while enabling feature reuse across teams.

\subsection{The Feature Store Problem}

Without a feature store, organizations face:

\begin{itemize}
    \item \textbf{Training-Serving Skew:} Features computed differently in training vs serving
    \item \textbf{Feature Duplication:} Teams reimplementing the same features
    \item \textbf{No Feature Discovery:} Data scientists unaware of existing features
    \item \textbf{Inconsistent Definitions:} Same feature computed differently across teams
    \item \textbf{Point-in-Time Correctness:} Historical features contain data leakage
\end{itemize}

\subsection{Feature Store Architecture}

\begin{figure}[h]
\centering
\begin{tikzpicture}[
    node distance=1.5cm,
    box/.style={rectangle, draw, fill=blue!20, text width=2.2cm, text centered, minimum height=0.9cm, font=\small},
    storage/.style={cylinder, draw, fill=green!20, text width=1.8cm, text centered, minimum height=0.8cm, aspect=0.3, font=\small},
    arrow/.style={->, >=stealth, thick}
]
    % Top: Feature definitions
    \node[box] (definitions) {Feature Definitions};

    % Middle: Ingestion
    \node[box, below left of=definitions, xshift=-2cm, yshift=-1.5cm] (batch) {Batch Pipeline};
    \node[box, below right of=definitions, xshift=2cm, yshift=-1.5cm] (stream) {Stream Pipeline};

    % Storage layers
    \node[storage, below of=definitions, yshift=-3.5cm] (offline) {Offline Store};
    \node[storage, right of=offline, xshift=3cm] (online) {Online Store};

    % Consumption
    \node[box, below of=offline, yshift=-1.5cm] (training) {Training};
    \node[box, below of=online, yshift=-1.5cm] (serving) {Serving};

    % Arrows
    \draw[arrow] (definitions) -- (batch);
    \draw[arrow] (definitions) -- (stream);
    \draw[arrow] (batch) -- (offline);
    \draw[arrow] (stream) -- (online);
    \draw[arrow] (offline) -- (training);
    \draw[arrow] (online) -- (serving);
    \draw[arrow, dashed] (offline) -- node[above, font=\tiny] {Materialize} (online);

    % Labels
    \node[left, text width=2cm, font=\tiny] at (offline.west) {Historical data, Parquet};
    \node[right, text width=2cm, font=\tiny] at (online.east) {Low latency, Redis};
\end{tikzpicture}
\caption{Feature Store Architecture}
\label{fig:feature_store}
\end{figure}

\subsection{Implementation with Feast}

Feast is an open-source feature store that supports both batch and real-time features:

\begin{lstlisting}[style=python, caption=Feature Store Implementation with Feast]
from feast import FeatureStore, Entity, FeatureView, Field
from feast.types import Float32, Int64, String
from feast.data_source import FileSource, PushSource
from datetime import timedelta
import pandas as pd
from typing import List, Dict

# Define feature repository structure
# features/repository.py

# Define entities
user = Entity(
    name="user_id",
    description="User identifier",
    join_keys=["user_id"]
)

# Batch feature source
user_activity_source = FileSource(
    path="data/user_activity.parquet",
    timestamp_field="event_timestamp",
    created_timestamp_column="created_timestamp"
)

# Define feature views
user_activity_features = FeatureView(
    name="user_activity_features",
    entities=[user],
    ttl=timedelta(days=30),
    schema=[
        Field(name="total_purchases", dtype=Int64),
        Field(name="avg_purchase_value", dtype=Float32),
        Field(name="days_since_last_purchase", dtype=Int64),
        Field(name="total_page_views", dtype=Int64),
        Field(name="session_count_7d", dtype=Int64),
        Field(name="cart_abandonment_rate", dtype=Float32)
    ],
    source=user_activity_source,
    online=True,
    tags={"team": "ml-platform", "domain": "user-behavior"}
)

# Real-time feature source
user_realtime_source = PushSource(
    name="user_realtime_features",
    batch_source=user_activity_source
)

user_realtime_features = FeatureView(
    name="user_realtime_features",
    entities=[user],
    ttl=timedelta(hours=1),
    schema=[
        Field(name="session_duration", dtype=Float32),
        Field(name="pages_viewed_session", dtype=Int64),
        Field(name="current_cart_value", dtype=Float32)
    ],
    source=user_realtime_source,
    online=True
)

class FeatureRegistry:
    """Centralized feature management"""

    def __init__(self, repo_path: str):
        self.store = FeatureStore(repo_path=repo_path)

    def get_historical_features(
        self,
        entity_df: pd.DataFrame,
        features: List[str]
    ) -> pd.DataFrame:
        """Get point-in-time correct historical features"""

        # entity_df must contain entity keys and event_timestamp
        training_df = self.store.get_historical_features(
            entity_df=entity_df,
            features=features
        ).to_df()

        return training_df

    def get_online_features(
        self,
        entity_rows: List[Dict[str, any]],
        features: List[str]
    ) -> Dict:
        """Get real-time features for prediction"""

        feature_vector = self.store.get_online_features(
            features=features,
            entity_rows=entity_rows
        ).to_dict()

        return feature_vector

    def materialize_features(
        self,
        start_date: str,
        end_date: str,
        feature_views: List[str] = None
    ):
        """Materialize batch features to online store"""

        self.store.materialize(
            start_date=start_date,
            end_date=end_date,
            feature_views=feature_views
        )

    def push_realtime_features(
        self,
        push_source_name: str,
        df: pd.DataFrame
    ):
        """Push real-time features to online store"""

        self.store.push(
            push_source_name=push_source_name,
            df=df
        )

# Usage for Training
def get_training_data():
    """Get features for training with point-in-time correctness"""

    registry = FeatureRegistry(repo_path="features/")

    # Define entity DataFrame with timestamps
    entity_df = pd.DataFrame({
        "user_id": ["user_1", "user_2", "user_3"],
        "event_timestamp": [
            pd.Timestamp("2024-01-15 10:00:00"),
            pd.Timestamp("2024-01-15 11:00:00"),
            pd.Timestamp("2024-01-15 12:00:00")
        ],
        "label": [1, 0, 1]  # Target variable
    })

    # Request features
    features = [
        "user_activity_features:total_purchases",
        "user_activity_features:avg_purchase_value",
        "user_activity_features:days_since_last_purchase",
        "user_activity_features:session_count_7d"
    ]

    # Get historical features (point-in-time correct)
    training_df = registry.get_historical_features(
        entity_df=entity_df,
        features=features
    )

    return training_df

# Usage for Serving
def get_serving_features(user_id: str) -> Dict:
    """Get features for real-time prediction"""

    registry = FeatureRegistry(repo_path="features/")

    # Request features for a single user
    features = [
        "user_activity_features:total_purchases",
        "user_activity_features:avg_purchase_value",
        "user_realtime_features:session_duration",
        "user_realtime_features:current_cart_value"
    ]

    feature_vector = registry.get_online_features(
        entity_rows=[{"user_id": user_id}],
        features=features
    )

    return feature_vector
\end{lstlisting}

\subsection{Feature Versioning and Lineage}

Feature stores enable tracking feature evolution and lineage:

\begin{lstlisting}[style=python, caption=Feature Versioning and Lineage]
from dataclasses import dataclass
from datetime import datetime
from typing import Optional, List
import hashlib
import json

@dataclass
class FeatureVersion:
    """Track feature versions and lineage"""

    name: str
    version: str
    definition: str
    created_at: datetime
    created_by: str
    source_tables: List[str]
    transformation_hash: str
    parent_version: Optional[str] = None

    @classmethod
    def from_transformation(
        cls,
        name: str,
        definition: str,
        source_tables: List[str],
        created_by: str
    ):
        """Create versioned feature from transformation"""

        # Hash the transformation logic for version control
        transformation_str = json.dumps({
            "definition": definition,
            "sources": sorted(source_tables)
        }, sort_keys=True)

        transformation_hash = hashlib.sha256(
            transformation_str.encode()
        ).hexdigest()[:8]

        # Version format: v{YYYYMMDD}_{hash}
        version = f"v{datetime.now().strftime('%Y%m%d')}_{transformation_hash}"

        return cls(
            name=name,
            version=version,
            definition=definition,
            created_at=datetime.now(),
            created_by=created_by,
            source_tables=source_tables,
            transformation_hash=transformation_hash
        )

# Example: Version a feature
feature = FeatureVersion.from_transformation(
    name="user_purchase_velocity",
    definition="""
        SELECT
            user_id,
            COUNT(*) / NULLIF(MAX(days_since_first_purchase), 0) as purchase_velocity
        FROM purchases
        GROUP BY user_id
    """,
    source_tables=["purchases"],
    created_by="data_scientist_1"
)

print(f"Feature {feature.name} version {feature.version}")
# Output: Feature user_purchase_velocity version v20240115_a3f2e9b1
\end{lstlisting}

\section{Schema Evolution and Compatibility}

As ML systems evolve, data schemas must change while maintaining backward compatibility.

\subsection{Schema Evolution Strategies}

\begin{table}[h]
\centering
\caption{Schema Evolution Strategies}
\label{tab:schema_evolution}
\begin{tabular}{@{}lp{5cm}p{5cm}@{}}
\toprule
\textbf{Strategy} & \textbf{Use Case} & \textbf{Compatibility} \\
\midrule
Add optional field & New feature without breaking old code & Backward compatible \\
Remove field & Deprecated feature & Forward compatible only \\
Rename field & Clarify semantics & Breaking change \\
Change type & Fix data type issue & Breaking change \\
Add default value & Make field required & Backward compatible \\
\bottomrule
\end{tabular}
\end{table}

\subsection{Schema Management with Avro}

Apache Avro provides schema evolution with compatibility checking:

\begin{lstlisting}[style=python, caption=Schema Evolution with Avro]
from dataclasses import dataclass
from typing import Optional
import avro.schema
from avro.io import DatumWriter, DatumReader
from io import BytesIO
import json

# Schema Version 1: Initial schema
schema_v1 = {
    "type": "record",
    "name": "UserEvent",
    "namespace": "ml.features",
    "fields": [
        {"name": "user_id", "type": "string"},
        {"name": "event_type", "type": "string"},
        {"name": "timestamp", "type": "long"},
        {"name": "value", "type": "double"}
    ]
}

# Schema Version 2: Add optional field (backward compatible)
schema_v2 = {
    "type": "record",
    "name": "UserEvent",
    "namespace": "ml.features",
    "fields": [
        {"name": "user_id", "type": "string"},
        {"name": "event_type", "type": "string"},
        {"name": "timestamp", "type": "long"},
        {"name": "value", "type": "double"},
        {"name": "session_id", "type": ["null", "string"], "default": None}
    ]
}

# Schema Version 3: Add field with default (backward compatible)
schema_v3 = {
    "type": "record",
    "name": "UserEvent",
    "namespace": "ml.features",
    "fields": [
        {"name": "user_id", "type": "string"},
        {"name": "event_type", "type": "string"},
        {"name": "timestamp", "type": "long"},
        {"name": "value", "type": "double"},
        {"name": "session_id", "type": ["null", "string"], "default": None},
        {"name": "device_type", "type": "string", "default": "unknown"}
    ]
}

class SchemaRegistry:
    """Manage schema versions and compatibility"""

    def __init__(self):
        self.schemas = {}

    def register_schema(self, name: str, version: int, schema_json: dict):
        """Register a new schema version"""
        schema = avro.schema.parse(json.dumps(schema_json))

        # Check compatibility with previous version
        if version > 1:
            previous_schema = self.schemas[(name, version - 1)]
            if not self._is_backward_compatible(previous_schema, schema):
                raise ValueError(
                    f"Schema {name} v{version} is not backward compatible"
                )

        self.schemas[(name, version)] = schema

    def _is_backward_compatible(
        self,
        reader_schema: avro.schema.Schema,
        writer_schema: avro.schema.Schema
    ) -> bool:
        """Check if writer schema is backward compatible with reader"""
        # Simplified check - production would use Avro's compatibility rules
        # New schema can read data written with old schema
        return True

    def get_schema(self, name: str, version: int) -> avro.schema.Schema:
        """Retrieve schema by name and version"""
        return self.schemas.get((name, version))

# Example usage
registry = SchemaRegistry()

registry.register_schema("UserEvent", 1, schema_v1)
registry.register_schema("UserEvent", 2, schema_v2)
registry.register_schema("UserEvent", 3, schema_v3)
\end{lstlisting}

\section{Data Versioning with DVC}

Data Version Control (DVC) enables versioning large datasets alongside code, ensuring reproducibility.

\subsection{DVC Workflow}

\begin{lstlisting}[style=bash, caption=DVC Workflow for Data Versioning]
# Initialize DVC in project
dvc init

# Configure remote storage (S3, GCS, Azure, etc.)
dvc remote add -d storage s3://my-ml-data/dvc-store

# Track data directory
dvc add data/raw/train.csv
# Creates data/raw/train.csv.dvc (metadata file)
# Adds data/raw/train.csv to .gitignore

# Commit DVC metadata to git
git add data/raw/train.csv.dvc .gitignore
git commit -m "Add training data v1"

# Push data to remote storage
dvc push

# Tag this version
git tag -a "data-v1.0" -m "Initial training data"
git push origin data-v1.0

# Later: Update dataset
# Modify data/raw/train.csv
dvc add data/raw/train.csv
git add data/raw/train.csv.dvc
git commit -m "Update training data with new samples"
git tag -a "data-v1.1" -m "Added 10k new samples"
dvc push
git push origin data-v1.1

# Checkout specific data version
git checkout data-v1.0
dvc checkout  # Downloads data-v1.0 from remote

# Create pipeline that tracks data dependencies
dvc run -n preprocess \
    -d data/raw/train.csv \
    -o data/processed/train_processed.parquet \
    python scripts/preprocess.py

# Pipeline automatically tracks lineage
dvc dag  # Show pipeline DAG
\end{lstlisting}

\subsection{DVC for Experiment Reproducibility}

\begin{lstlisting}[style=python, caption=Reproducible Experiments with DVC]
# dvc.yaml - Define ML pipeline
stages:
  prepare:
    cmd: python src/prepare.py
    deps:
      - src/prepare.py
      - data/raw
    params:
      - prepare.test_split
      - prepare.seed
    outs:
      - data/prepared

  featurize:
    cmd: python src/featurize.py
    deps:
      - src/featurize.py
      - data/prepared
    params:
      - featurize.max_features
      - featurize.ngrams
    outs:
      - data/features

  train:
    cmd: python src/train.py
    deps:
      - src/train.py
      - data/features
    params:
      - train.n_estimators
      - train.max_depth
    outs:
      - models/model.pkl
    metrics:
      - metrics/train_metrics.json:
          cache: false

# params.yaml - Pipeline parameters
prepare:
  test_split: 0.2
  seed: 42

featurize:
  max_features: 500
  ngrams: 2

train:
  n_estimators: 100
  max_depth: 10

# Run entire pipeline
# dvc repro

# Compare experiments
# dvc params diff
# dvc metrics diff
\end{lstlisting}

\section{Complete Production Data Pipeline Example}

Integrating all concepts into a production-grade pipeline:

\begin{lstlisting}[style=python, caption=End-to-End Production Data Pipeline]
from typing import Dict, Any
import logging
from datetime import datetime
from airflow import DAG
from airflow.operators.python import PythonOperator
from airflow.providers.google.cloud.operators.bigquery import BigQueryInsertJobOperator
import great_expectations as ge

logger = logging.getLogger(__name__)

# Define DAG
default_args = {
    'owner': 'ml-team',
    'depends_on_past': False,
    'email_on_failure': True,
    'email': ['ml-alerts@company.com'],
    'retries': 2
}

dag = DAG(
    'ml_data_pipeline',
    default_args=default_args,
    description='Production ML data pipeline',
    schedule_interval='0 2 * * *',  # Daily at 2 AM
    start_date=datetime(2024, 1, 1),
    catchup=False,
    tags=['ml', 'production']
)

def extract_data(**context):
    """Extract data from source systems"""
    from google.cloud import bigquery

    client = bigquery.Client()

    # Extract data with incremental loading
    execution_date = context['ds']

    query = f"""
        SELECT *
        FROM `project.dataset.events`
        WHERE DATE(event_timestamp) = '{execution_date}'
    """

    df = client.query(query).to_dataframe()

    # Save to intermediate storage
    output_path = f"gs://data-lake/raw/events/{execution_date}.parquet"
    df.to_parquet(output_path)

    return output_path

def validate_data(**context):
    """Validate data quality"""
    from ml_validator import MLDataValidator
    import pandas as pd

    # Get path from previous task
    data_path = context['ti'].xcom_pull(task_ids='extract')

    # Load data
    df = pd.read_parquet(data_path)

    # Run validation
    validator = MLDataValidator(context_root_dir='./ge_config')

    results = validator.validate_data(
        df=df,
        suite_name='events_schema',
        run_name=context['ds']
    )

    # Fail task if validation fails
    if not results['success']:
        raise ValueError(
            f"Validation failed: {results['unsuccessful_expectations']} "
            f"expectations failed"
        )

    # Check for drift
    reference_df = pd.read_parquet('gs://data-lake/reference/events.parquet')
    drift_results = validator.detect_data_drift(reference_df, df)

    drifted_features = sum(1 for r in drift_results.values() if r['drift_detected'])

    if drifted_features > 5:
        logger.warning(
            f"High drift detected: {drifted_features} features drifted"
        )

    return data_path

def transform_features(**context):
    """Transform and engineer features"""
    import pandas as pd
    from feature_engineering import FeatureEngineer

    data_path = context['ti'].xcom_pull(task_ids='validate')
    df = pd.read_parquet(data_path)

    # Apply feature engineering
    engineer = FeatureEngineer()
    features_df = engineer.transform(df)

    # Save features
    output_path = f"gs://data-lake/features/{context['ds']}.parquet"
    features_df.to_parquet(output_path)

    return output_path

def load_feature_store(**context):
    """Load features into feature store"""
    from feature_store import FeatureRegistry
    import pandas as pd

    features_path = context['ti'].xcom_pull(task_ids='transform')
    features_df = pd.read_parquet(features_path)

    # Push to feature store
    registry = FeatureRegistry(repo_path='feature_repo/')

    # Push to offline store (for training)
    features_df.to_parquet('data/features/latest.parquet')

    # Materialize to online store (for serving)
    registry.materialize_features(
        start_date=context['ds'],
        end_date=context['ds'],
        feature_views=['user_features']
    )

    logger.info(f"Loaded {len(features_df)} feature rows to feature store")

# Define tasks
extract_task = PythonOperator(
    task_id='extract',
    python_callable=extract_data,
    dag=dag
)

validate_task = PythonOperator(
    task_id='validate',
    python_callable=validate_data,
    dag=dag
)

transform_task = PythonOperator(
    task_id='transform',
    python_callable=transform_features,
    dag=dag
)

load_task = PythonOperator(
    task_id='load_feature_store',
    python_callable=load_feature_store,
    dag=dag
)

# Set dependencies
extract_task >> validate_task >> transform_task >> load_task
\end{lstlisting}

\section{Chapter Summary}

This chapter covered the complete data pipeline architecture for ML systems:

\begin{itemize}
    \item \textbf{Data Ingestion:} Batch, streaming, and hybrid patterns using Apache Beam and Kafka
    \item \textbf{ETL vs ELT:} Modern ELT approach preserves raw data for ML flexibility
    \item \textbf{Data Validation:} Comprehensive validation with Great Expectations and drift detection
    \item \textbf{Feature Stores:} Centralized feature management with Feast for consistency
    \item \textbf{Schema Evolution:} Managing schema changes with Avro compatibility rules
    \item \textbf{Data Versioning:} Reproducible datasets with DVC
    \item \textbf{Production Integration:} End-to-end pipelines with Airflow orchestration
\end{itemize}

Key takeaways:

\begin{enumerate}
    \item Data pipelines are the foundation of ML systems—invest heavily in quality and reliability
    \item Choose ELT over ETL for ML to preserve raw data and enable experimentation
    \item Implement comprehensive validation to catch data quality issues early
    \item Use feature stores to eliminate training-serving skew and enable feature reuse
    \item Version data alongside code for reproducibility
    \item Monitor data drift continuously—models degrade when data distributions change
    \item Design for failure: dead letter queues, retry logic, and alerting
\end{enumerate}

\section{Further Reading}

\begin{itemize}
    \item Kleppmann, M. (2017). \textit{Designing Data-Intensive Applications}. O'Reilly.
    \item Reis, J., \& Housley, M. (2022). \textit{Fundamentals of Data Engineering}. O'Reilly.
    \item Apache Beam Programming Guide: \url{https://beam.apache.org/documentation/}
    \item Great Expectations Documentation: \url{https://docs.greatexpectations.io/}
    \item Feast Feature Store: \url{https://feast.dev/}
    \item DVC Documentation: \url{https://dvc.org/doc}
\end{itemize}

\section{Exercises}

\begin{exercise}
\textbf{Batch Pipeline Implementation}

Implement a complete batch data pipeline:

\begin{enumerate}
    \item Create an Apache Beam pipeline that:
    \begin{itemize}
        \item Reads CSV files from GCS/S3
        \item Validates data schema and ranges
        \item Handles invalid records with a dead letter queue
        \item Aggregates user-level metrics
        \item Writes results to BigQuery
    \end{itemize}
    \item Add error handling and retry logic
    \item Implement Dataflow runner configuration
    \item Add monitoring with custom metrics
\end{enumerate}
\end{exercise}

\begin{exercise}
\textbf{Streaming Pipeline with Kafka}

Build a real-time streaming pipeline:

\begin{enumerate}
    \item Set up Kafka topics for event ingestion
    \item Implement a consumer that:
    \begin{itemize}
        \item Processes events in real-time
        \item Maintains sliding window aggregations
        \item Handles late-arriving data
        \item Implements exactly-once semantics
    \end{itemize}
    \item Write processed features to a feature store
    \item Add windowed aggregations (5-min, 1-hour windows)
    \item Implement state management for user-level features
\end{enumerate}
\end{exercise}

\begin{exercise}
\textbf{Data Validation Suite}

Create a comprehensive validation framework:

\begin{enumerate}
    \item Initialize Great Expectations project
    \item Create expectation suites for:
    \begin{itemize}
        \item Schema validation
        \item Statistical property checks
        \item Cross-column relationships
        \item Custom business logic
    \end{itemize}
    \item Implement drift detection using KS tests
    \item Generate data quality reports
    \item Integrate validation into CI/CD pipeline
    \item Set up alerting for validation failures
\end{enumerate}
\end{exercise}

\begin{exercise}
\textbf{Feature Store Implementation}

Build a feature store with Feast:

\begin{enumerate}
    \item Define feature repository structure
    \item Create batch and streaming feature views
    \item Implement offline store (Parquet/BigQuery)
    \item Configure online store (Redis/DynamoDB)
    \item Materialize batch features to online store
    \item Implement point-in-time correct feature retrieval for training
    \item Create serving endpoint for real-time predictions
    \item Add feature versioning and lineage tracking
\end{enumerate}
\end{exercise}

\begin{exercise}
\textbf{Schema Evolution Strategy}

Design schema evolution framework:

\begin{enumerate}
    \item Define Avro schemas for event data
    \item Implement schema registry with version control
    \item Add backward-compatible schema changes:
    \begin{itemize}
        \item New optional field
        \item Field with default value
        \item Deprecated field handling
    \end{itemize}
    \item Write compatibility tests
    \item Document migration strategy
    \item Implement gradual rollout process
\end{enumerate}
\end{exercise}

\begin{exercise}
\textbf{Data Versioning with DVC}

Implement data version control:

\begin{enumerate}
    \item Initialize DVC project
    \item Configure remote storage (S3/GCS)
    \item Version training and test datasets
    \item Create DVC pipeline with stages:
    \begin{itemize}
        \item Data preparation
        \item Feature engineering
        \item Model training
    \end{itemize}
    \item Track parameter changes in params.yaml
    \item Compare experiments using dvc metrics
    \item Reproduce historical experiments
    \item Document data lineage
\end{enumerate}
\end{exercise}

\begin{exercise}
\textbf{End-to-End Production Pipeline}

Integrate all concepts into a production pipeline:

\begin{enumerate}
    \item Design Airflow DAG with tasks:
    \begin{itemize}
        \item Data extraction (batch/streaming)
        \item Validation with Great Expectations
        \item Feature engineering
        \item Feature store loading
        \item Data quality monitoring
    \end{itemize}
    \item Implement error handling and retries
    \item Add data quality gates (fail on validation errors)
    \item Set up drift detection and alerting
    \item Configure monitoring dashboards
    \item Implement data SLAs (freshness, completeness)
    \item Write integration tests
    \item Document pipeline operations
\end{enumerate}
\end{exercise}
